\setcounter{subsection}{1}
\subsection{Dualidade}

\begin{definition}
    Seja \(V\) \(\F\)-espaço vetorial. Denotamos por \(V':= \hom(V,\F)\), o \emph{espaço dual} de \(V\), \(f \in V'\) é chamado de \emph{funcional linear} em \(V\).   
\end{definition}

\begin{definition}
    Sejam \(W\) e \(V\) \(\F\)-espaços vetoriais. Chamamos de \emph{pareamiento bilinear} os elementos \(\phi \in \hom^2(W\times V,\F)\).  
\end{definition}

\begin{example}
    \(\operatorname{av}: V' \times V \ni (f,v) \mapsto f(v) \in \F\). \comentario{Função avaliação, mostre que ta bém definida, ou seja, \(\operatorname{av} \in \hom^2(V'\times V, \F )\) }
\end{example}   
 
\begin{note}
    Segue da \(\square\) \ref{prop:9.1.4} que se \(\alpha = (v_i)\) é base de \(V\), então \(\alpha' = (f_i)\), onde \(f_i(v_j) = \delta_{ij}\) é l.i. em \(V'\), base no caso \(\dim V <\infty\). \footnote{Aquí pensamos em \(W=\F = [1]\), espaço vetorial de dimensão um sobre se mesmo. }%\(\square\) \ref{prop:9.1.4} segue que 
\end{note}

\ejemplo{Conta con duais \centering
\tcblower
No casso finito é fazer conta com matrices. De fato, se \(\dim V = n < \infty\) e \(\alpha =(v_j)\), \(\beta = (f_j)\) bases pra \(V\) e \(V'\), então  
\[\beta = \alpha ' \ \leftrightharpoons \ \begin{bmatrix}
   \rotatebox{90}{$|$} & f_1 & \rotatebox{90}{$|$} \\ 
    & \rotatebox{90}{$\ldots$}& \\ 
   \rotatebox{90}{$|$} & f_n & \rotatebox{90}{$|$}  
\end{bmatrix} \begin{bmatrix}
   \rotatebox{0}{$|$} &  & \rotatebox{0}{$|$} \\ 
   v_1 & \rotatebox{0}{$\cdots$}& v_n\\ 
   \rotatebox{0}{$|$} &  & \rotatebox{0}{$|$}  
\end{bmatrix} = \id_n.\]
Note também que \(\alpha \) e \(\beta \) são bases \(\leftrightharpoons\) as matrices acima são invertíveis.
\newline \vspace{-0.4cm}
\begin{example}
    Sejam \(V = \mathcal{P}_2(\R)\) e \(\beta = f_1, f_2, f_3\) dadas por, 
    \[\textstyle f_1(p) = p(1), \hspace{1cm} f_2(p) = \dot{p}(0),\hspace{1cm} f_3(p) = \int_0^{\sqrt{2}} \dot{p}(t)t^3 \ dt. \]
    Sendo \(p(t) = a + bt + ct^2 \in \mathcal{P}_2(\R)\) qualquer, temos \(\textstyle f_1(p) = a+b+c,\  f_2(p) = b,\ f_3(p) = b + \frac{8\sqrt{2}}{5} c= b+\lambda c\). Assim, em termos matriciais, 
    \[\begin{bmatrix}
        1 & 1 & 1\\ 
        0 & 1 & 0\\ 
        0 & 1 & \lambda\\ 
    \end{bmatrix}\begin{bmatrix}
        1 & \frac{1}{\lambda} -1 & -\frac{1}{\lambda}\\
        0 & 1 & 0\\
        0 & -\frac{1}{\lambda}  & \frac{1}{\lambda}\\
    \end{bmatrix} = \id_3.\]
    Visto que tudo esta pensado em termos da base canônica (\(1, t, t^2\)), temos 
    \[\textstyle\alpha = \left\{1, (\frac{1}{\lambda}-1) + t- \frac{1}{\lambda}t^2, -\frac{1}{\lambda} + \frac{1}{\lambda}t^2\right\} : \alpha' = \beta. \]
\end{example}
}

\hypertarget{corol921}{}
\begin{corollary}
    \textbf{9.2.1.} Seja \(v\in V\) tal que \(\forall f \in V',\ f(v) = 0 \), então \(v = 0\). \comentario{Se \(v\neq 0\), então \(\exists \alpha \) base de \(V\) tal que \(v = v_{i_0}\in \alpha\), assim \(f_{i_0}(v_{i_0}) = 1 \neq 0\) }
\end{corollary}
 
\propositionnum{9.2.2.}
\begin{proposition}
    \label{prop:9.2.2}
   Seja \(I\) tal que  \(|I| = \dim V\), então \(V'\cong \mathcal{F}(I, \F)\). 
\end{proposition}

\demo{
Sejam \(\alpha = (v_i)\) base de \(V\) e \(\mathcal{F} (I,\F) \ni (a_i) \sim a: i \mapsto a_i \in \F\). Trabalha com \(
    T: V' \ni f\mapsto (f(v_i)) \in \mathcal{F} (I,\F)\) e \(S: \mathcal{F} (I, \F) \ni (a_i) \mapsto (f: v_j \mapsto a_j)\), mostra que são lineares e inversas a uma da outra. 
}

\begin{corollary}
    Segue de \(\square\) \ref{prop:9.2.2} que se \(\dim V = \infty \), então \([\alpha'] \subset V'\). \comentario{Na demostração \(T([\alpha']) = \{(a_i) \in \mathcal{F}(I,\F) : a_i \neq 0 \text{ em finitos }i \} \subset \mathcal{F}(I, \F) \cong V'\)} 
\end{corollary}

\begin{note}
    Isso só mostra que \(\alpha'\) não é base de \(V'\), afirmar \(\dim V' > \dim V\) mesmo sendo certa, precisa mais coisa (analisé de cardinais).   
\end{note}

\propositionnum{}
\vspace{-0.5cm}
\begin{lemma}
    Sendo \(V\) \(\F\)-espaço vetorial a função \(J_{_V}: V \ni \textcolor{red}{v \mapsto} \overset{\underset{\rotatebox{90}{$\in$}}{V''}}{\textcolor{red}{J_{_V}(v)}}: \overset{\underset{\rotatebox{90}{$\in$}}{V'}}{f}\mapsto f(v) \in \F\) 
    é bém definida, linear e injetora.  
\end{lemma}

\demo{
    Bém definida é ver que \(J_{_V}(v) \in V''\), olha pra \(J_{_V}(v)(f+\lambda g)\). Linearidade é quase-imediata e a última afirmação segue do {\scriptsize\(\boxtimes\)} \hyperlink{corol921}{9.2.1.}.   
}

\begin{corollary}
    Se \(\dim V < \infty\), então \(\dim V = \dim V' = \dim V''\) e \(J_{_V}: V \cong V''\). 
\end{corollary}

\begin{lemma}
    Sejam \(V, \ W\) \(\F\)-espaços vetoriais e \(T\in \hom(V,W)\). O \emph{operador adjunto}, \(T': W' \ni f \mapsto f \circ T \in V' \in \hom(W',V')\). \comentario{Propiedades da composição} 
\end{lemma}

\propositionnum{9.2.3.}
\begin{proposition}
    \label{prop:9.2.3}
    Sejam \(V, W\) \(\F\)-espaços vetorias de dimensão finita, \(\alpha\), \(\beta \) bases pra eles respetivamente e \(T\in \hom(V,W)\). Então, \([T']^{\beta'}_{\alpha'} = ([T]^\alpha_\beta)^t\). 
\end{proposition}

\demo{
    Sejam \(\alpha = (v_j)\), \(\beta= (w_i)\) e  \([T]^\alpha_\beta = (a_{ij})\). Escreva \(T(v_j) = \sum^{m} a_{ij} w_i\), logo note que a entrada \(b_{ji}\) da \([T']^{\beta'}_{\alpha'}\) é precisamente \(T'(g_i)(v_j)\). 
}


